\documentclass[a4j,dvipdfmx]{jsarticle}
\usepackage{amsmath,amssymb}
\usepackage{siunitx}
\usepackage{ascmac}
\usepackage{amsthm}
\usepackage{bm}

\usepackage[margin=10truemm]{geometry}

% \usepackage[dvipdfmx]{hyperref}
\begin{document}
    \part*{第一回 線型代数 問題}
    \subsection*{問題1}
        次のベクトルが一次独立か一次従属かを判定する.
        \begin{equation*}
            \begin{pmatrix}1 \\ 2 \\ 3 \end{pmatrix}, \begin{pmatrix}4 \\ 5\\ 6\end{pmatrix}, \begin{pmatrix}7 \\ 8 \\ 9 \end{pmatrix}
        \end{equation*}
    \subsection*{問題2}
        次の連立一次方程式をとく.
        \begin{equation*}
            (1)\left\{\begin{array}{c}
                2x+3y+2z=4\\
                4x+3y+2z=5\\
                x+2y+2z=3
            \end{array}\right.\quad 
            (2)\left\{\begin{array}{c}
                x+2y+3z-4w=0\\
                x+y+z-w=0\\
                x+3y+5z-7w=0
            \end{array}\right.
        \end{equation*}
    \subsection*{問題3}
        次の行列式を計算する.
        \begin{equation*}
            (1)\begin{vmatrix}
                2 & 4 & 3 \\
                3 & 8 & 2\\
                2 & 8 & 6
            \end{vmatrix}\quad
            (2)\begin{vmatrix}
                1 & 2 & 3 & 4\\
                2 & 2 & 1 & 5\\
                3 & 3 & 2 & 8\\
                4 & 5 & 4 & 10
            \end{vmatrix}\quad
            (3)\begin{vmatrix}
                1 & 100 & 65 & 98765421 & 15\\
                0 & 8 & 23 & 15 & 31\\
                0 & 0 & 314 & 2 & 99\\
                0 & 0 & 0 & 0.625 & -1\\
                0 & 0 & 0 & 0 &  5
            \end{vmatrix}
        \end{equation*}
    \subsection*{問題4}
        次の行列の積を計算する. ただし, $\bm{a},\bm{b},\bm{c}\in \mathbb{R}^3$
        \begin{equation*}
            (1)\begin{pmatrix}
                1 & -2 & 3\\
                2 & 1 & 0
            \end{pmatrix}
            \begin{pmatrix}
                -2 & -1 & 1\\
                1 & 2 & -3\\
                4 & 0 & 3
            \end{pmatrix}\quad 
            (2)\begin{pmatrix}
                1 & -1 & 0 & 0\\
                0 & -2 & 0 & 0\\
                0 & 0  & -2 & 3\\
                0 & 0  & 1 & 1
            \end{pmatrix}
            \begin{pmatrix}
                2 & 1 & 0 & 0\\
                0 & 1 & 0 & 0\\
                0 & 0  & 1 & 1\\
                0 & 0  & 2 & -3
            \end{pmatrix}\quad
            (3)\begin{pmatrix}
                \bm{a} & \bm{b} & \bm{c}
            \end{pmatrix}
            \begin{pmatrix}
                x\\
                y\\
                z
            \end{pmatrix}
        \end{equation*}
    \subsection*{問題5}
        逆行列の掃き出し法による求め方の原理を説明せよ.
    \subsection*{問題6}
        次の命題を示せ.
        \begin{equation*}
            (1)A\subset A \quad (2)A\cap B\subset A\quad (3)A\cap A = A
        \end{equation*}
    \subsection*{問題7}
        $f:\mathbb{R}\to\mathbb{R},x\mapsto \sin x$について, 次を求めよ. 
        \begin{equation*}
            (1)f\left([-\frac{\pi}{3},\frac{\pi}{3}]\right)\quad (2)f\left([-\frac{\pi}{2},\frac{\pi}{2}]\right)\quad (3)f\left([-2\pi,\pi]\right)\quad (4)f^{-1}([-2,0])
        \end{equation*}
    \subsection*{問題8}
        写像$f:\mathbb{R}\to\mathbb{R}, x\mapsto ax+b\hspace{1mm}(a\neq 0)$は全単射であることを(定義より)示せ.
    \subsection*{問題9}
        写像$f:X\to Y,g:Y\to Z$がどちらも全単射であるとき, 合成写像$g\circ f:X\to Z$も全単射であることを(定義より)示せ.
\end{document}